\subsection{QUBO Formulations}\label{sec:qubo-formulations}

In this section we present a few known QUBO formulations.

\begin{flushleft}
    \textcolor{Green3}{\faIcon{question-circle} \textbf{What is a QUBO formulation?}}
\end{flushleft}
A \definition{Quadratic Unconstrained Binary Optimization (QUBO)} formulation \hl{rewrites an optimization problem as the minimization of a quadratic function of binary variables:}
\begin{equation*}
    \min_{\mathbf{x}\in\{0,1\}^{n}} E(\mathbf{x})
    =
    \min_{\mathbf{x}\in\{0,1\}^{n}}
    \mathbf{x}^{T}Q\mathbf{x}
\end{equation*}
The variables encode candidate decisions, while the energy function assigns lower
values to better solutions. Since the formulation is \textbf{unconstrained}, hard
constraints are typically converted into \textbf{penalty terms}: the penalty is zero when a constraint is satisfied and positive when it is violated. A coefficient such as $P$ or $p$ controls the strength of that penalty.

\highspace
The formulas below are useful as ready-made encodings. In practice, the usual
workflow is:
\begin{enumerate}
    \item \textbf{Choose binary variables} to encode the problem.
    \item \textbf{Write the objective function} in terms of those variables.
    \item \textbf{Add penalty terms} for any hard constraints.
    \item \textbf{Expand} the resulting expression into a \textbf{quadratic polynomial}.
\end{enumerate}
The coefficients of the resulting polynomial can then be placed in the QUBO
matrix $Q$.

\begin{flushleft}
    \textcolor{Green3}{\faIcon{book} \textbf{General Penalty Conversion Rules}}
\end{flushleft}

\begin{table}[!htp]
    \centering
    \begin{tabular}{@{} l l @{}}
        \toprule
        \textbf{Classical Constraint} & \textbf{Equivalent Penalty} \\
        \midrule
        $x+y\leq 1$
        & $P(xy)$ \\[.3em]

        $x+y\geq 1$
        & $P(1-x-y+xy)$ \\[.3em]

        $x+y=1$
        & $P(1-x-y+2xy)$ \\[.3em]

        $x\leq y$
        & $P(x-xy)$ \\[.3em]

        $x_1+x_2+x_3\leq 1$
        & $P(x_1x_2+x_1x_3+x_2x_3)$ \\[.3em]

        $x=y$
        & $P(x+y-2xy)$ \\
        \bottomrule
    \end{tabular}
    \caption{A few known constraint/penalty pairs.}
    \label{tab:qubo-penalty-conversion}
\end{table}

\noindent
The following notation is used in the formulas below:
\begin{itemize}
    \item $G=(V,E)$ is an undirected graph.

    \item In graph problems, assume $(i,j)=(j,i)$ and then consider each edge only once.

    \item The complement graph is
        \begin{equation*}
            \bar{G}=(V,\bar{E}),
            \qquad
            \bar{E}
            =
            \left\{
                (i,j) : i\neq j,\ (i,j)\notin E
            \right\}
        \end{equation*}

    \item In Set Packing, $c_i$ is the capacity of set $S_i$,
        $B\in\{0,1\}^{m\times n}$ is a constraint matrix
        ($m$ constraints and $n$ variables), and
        $\vec{1}$ is a vector of $n$ elements, all equal to $1$.

    \item In CSPs, $C$ is a symmetric cost matrix, $A$ is a constraint
        matrix, and $b$ is a target vector.
\end{itemize}

\longline

\subsubsection{Max-Cut}

\textcolor{Green3}{\faIcon{question-circle} \textbf{What it solves.}} Partition the nodes of a graph into two sets $S$ and $V-S$ so as to maximize the number of edges crossing between them.

\highspace
\textcolor{Green3}{\faIcon{book} \textbf{Variable definition.}}
\begin{equation*}
    x_i=
    \begin{cases}
        1, & i\in S,\\
        0, & \text{otherwise}.
    \end{cases}
\end{equation*}

\highspace
\textcolor{Red2}{\faIcon{dollar-sign} \textbf{Cost function.}}
\begin{equation}\label{eq:qubo-max-cut}
    \min
    -\sum_{(i,j)\in E}
    \left(x_i+x_j-2x_ix_j\right)
\end{equation}

\highspace
\textcolor{Red2}{\faIcon{exclamation-triangle} \textbf{Penalty.}} None.

\longline

\subsubsection{Maximum Independent Set}

\textcolor{Green3}{\faIcon{question-circle} \textbf{What it solves.}} Find the largest set of nodes $S$ with no edges connecting them.

\highspace
\textcolor{Green3}{\faIcon{book} \textbf{Variable definition.}}
\begin{equation*}
    x_i=
    \begin{cases}
        1, & i\in S,\\
        0, & \text{otherwise}.
    \end{cases}
\end{equation*}

\highspace
\textcolor{Red2}{\faIcon{dollar-sign} \textbf{Cost function.}}
\begin{equation}\label{eq:qubo-max-independent-set}
    \min -\sum_{i}^{|V|}x_i
\end{equation}

\highspace
\textcolor{Red2}{\faIcon{exclamation-triangle} \textbf{Penalty.}}
\begin{equation}\label{eq:qubo-max-independent-set-penalty}
    p\sum_{(i,j)\in E}x_ix_j,
    \qquad p=2
\end{equation}

\newpage

\subsubsection{Minimum Vertex Cover}

\textcolor{Green3}{\faIcon{question-circle} \textbf{What it solves.}} Find the smallest set of nodes $C$ that touches every edge in the graph.

\highspace
\textcolor{Green3}{\faIcon{book} \textbf{Variable definition.}}
\begin{equation*}
    x_i=
    \begin{cases}
        1, & i\in C,\\
        0, & \text{otherwise}.
    \end{cases}
\end{equation*}

\highspace
\textcolor{Red2}{\faIcon{dollar-sign} \textbf{Cost function.}}
\begin{equation}\label{eq:qubo-min-vertex-cover}
    \min \sum_{i}^{|V|}x_i
\end{equation}

\highspace
\textcolor{Red2}{\faIcon{exclamation-triangle} \textbf{Penalty.}}
\begin{equation}\label{eq:qubo-min-vertex-cover-penalty}
    p\sum_{(i,j)\in E}
    \left(1-x_i-x_j+x_ix_j\right),
    \qquad p=|V|
\end{equation}

\longline

\subsubsection{Maximum Clique}

\textcolor{Green3}{\faIcon{question-circle} \textbf{What it solves.}} Find the largest fully connected subset of nodes (clique) $C$.

\highspace
\textcolor{Green3}{\faIcon{book} \textbf{Variable definition.}}
\begin{equation*}
    x_i=
    \begin{cases}
        1, & i\in C,\\
        0, & \text{otherwise}.
    \end{cases}
\end{equation*}

\highspace
\textcolor{Red2}{\faIcon{dollar-sign} \textbf{Cost function.}} This is the Maximum Independent Set formulation applied to the complement graph $\bar{G}$:
\begin{equation}\label{eq:qubo-max-clique}
    \min -\sum_{i}^{|V|}x_i
\end{equation}

\highspace
\textcolor{Red2}{\faIcon{exclamation-triangle} \textbf{Penalty.}}
\begin{equation}\label{eq:qubo-max-clique-penalty}
    p\sum_{(i,j)\in\bar{E}}x_ix_j,
    \qquad p=2
\end{equation}

\newpage

\subsubsection{Number Partitioning}

\textcolor{Green3}{\faIcon{question-circle} \textbf{What it solves.}} Partition a set of numbers $Z$ into two groups $Z_1,Z_2$ such that the sums of the elements in $Z_1$ and $Z_2$ are as equal as possible.

\highspace
\textcolor{Green3}{\faIcon{book} \textbf{Variable definition.}}
\begin{equation*}
    x_i=
    \begin{cases}
        1, & z_i\in Z_1,\\
        0, & z_i\in Z_2.
    \end{cases}
\end{equation*}

\highspace
\textcolor{Red2}{\faIcon{dollar-sign} \textbf{Cost function.}}
\begin{equation}\label{eq:qubo-number-partitioning}
    \min
    \left(
        \sum_{i}^{|Z|}
        (2x_i-1)z_i
    \right)^2
\end{equation}

\highspace
\textcolor{Red2}{\faIcon{exclamation-triangle} \textbf{Penalty.}} None.

\longline

\subsubsection{2-SAT}\label{sec:qubo-2-sat}

\textcolor{Green3}{\faIcon{question-circle} \textbf{What it solves.}} Find a Boolean assignment that satisfies all clauses with two literals each.

\highspace
\textcolor{Green3}{\faIcon{book} \textbf{Variable definition.}} One binary variable $x_i$ per Boolean literal.

\highspace
\textcolor{Red2}{\faIcon{dollar-sign} \textbf{Cost function.}} Sum of the penalties associated with the individual clauses.

\highspace
\textcolor{Red2}{\faIcon{exclamation-triangle} \textbf{Clause penalties.}}
\begin{equation}\label{eq:qubo-2-sat-penalties}
    \begin{aligned}
        x_i\lor x_j
        &\Rightarrow 1-x_i-x_j+x_ix_j,\\
        x_i\lor \bar{x}_j
        &\Rightarrow x_j-x_ix_j,\\
        \bar{x}_i\lor \bar{x}_j
        &\Rightarrow x_ix_j
    \end{aligned}
\end{equation}

\newpage

\subsubsection{Set Packing}

\textcolor{Green3}{\faIcon{question-circle} \textbf{What it solves.}} Select the most valuable disjoint subset of sets under mutual exclusion constraints.

\highspace
\textcolor{Green3}{\faIcon{book} \textbf{Variable definition.}}
\begin{equation*}
    x_i=
    \begin{cases}
        1, & S_i\text{ is selected},\\
        0, & \text{otherwise}.
    \end{cases}
\end{equation*}

\highspace
\textcolor{Red2}{\faIcon{dollar-sign} \textbf{Cost function.}}
\begin{equation}\label{eq:qubo-set-packing}
    \min -\sum_{i}^{n}c_ix_i,\quad \text{subject to}\qquad Bx\leq\vec{1}
\end{equation}

\highspace
\textcolor{Red2}{\faIcon{exclamation-triangle} \textbf{Penalty.}} Use the first and fifth formulas in \autoref{tab:qubo-penalty-conversion}, with:
\begin{equation}\label{eq:qubo-set-packing-penalty}
    p=\sum_{i}^{n}c_i
\end{equation}

\longline

\subsubsection{Constraint Satisfaction Problem (CSP)}

\textcolor{Green3}{\faIcon{question-circle} \textbf{What it solves.}} Minimize a cost function while satisfying linear equality constraints.

\highspace
\textcolor{Green3}{\faIcon{book} \textbf{Variable definition.}}
\begin{equation*}
    x_i\in\{0,1\}
\end{equation*}

\highspace
\textcolor{Red2}{\faIcon{dollar-sign} \textbf{Cost function.}}
\begin{equation}\label{eq:qubo-csp-cost}
    \min x^{T}Cx,\quad \text{subject to}\quad Ax=b
\end{equation}

\highspace
\textcolor{Red2}{\faIcon{exclamation-triangle} \textbf{Penalty.}}
\begin{equation}\label{eq:qubo-csp-penalty}
    p(Ax-b)^{T}(Ax-b)
\end{equation}
